% ============================================================
%  07_validation.tex -- Validation and final comparison
% ============================================================

\chapter{Validation and Final Comparison}
\label{ch:validation}

Chapter~\ref{ch:internals} described how FeynPy turns a declared Lagrangian
into a Feynman-rule \texttt{Expression}. This chapter checks that output
against an independent reference: the FeynRules exports of the same models.
FeynRules is treated as the mathematical oracle throughout.

The two outputs are not expected to be \emph{identical} as printed
expressions. FeynRules and FeynPy package charge conjugation differently,
expand covariant derivatives to different depths, and represent chirality
differently (an explicit projector in FeynRules, a field class in FeynPy).
What is checked instead is that the two outputs are \emph{equal as tensors}:
that they differ only by index relabelling, reordering of commuting factors,
and a small, explicitly listed set of gauge and Lorentz identities. Section
\ref{sec:validation-canonical} makes this precise and shows, with two worked
examples, what that equality test actually does.

The headline results: the packaged Standard Model matches its FeynRules
export exactly ($163/163$ vertices), the unbroken background-field Standard
Model matches its own export exactly ($67/67$), and the SMEFT2 Green-basis
model matches all $184$ exported rows of the FeynRules EFT-only
\texttt{Ltot}. Of those $184$ SMEFT2 rows, $161$ match with no assumption at
all, $15$ match after one globally fixed charge-conjugation convention, and
$8$ match after an additional, explicitly tabulated row-specific transform.
No row is accepted by searching over signs or field orderings, and no row is
left unresolved.

% --------------------------------------------------------------
\section{Validation strategy}
\label{sec:validation-strategy}
% --------------------------------------------------------------

The validation has three layers, and they must not be conflated: the outer
two bound what the inner one can mean. A perfect equality score over a small,
cherry-picked subset of a model would prove very little.

\paragraph{Model fidelity.} Does the FeynPy source transcribe every operator
of the FeynRules source? A question about the \emph{input}, answered in
\Cref{sec:validation-fidelity} by a block-by-block and
coefficient-by-coefficient audit.

\paragraph{Vertex equality.} Do the Feynman rules derived from the two
Lagrangians agree as tensor expressions? The substance of the comparison,
answered in \Cref{sec:validation-canonical} and
\Cref{sec:validation-smeft2,sec:validation-ec}.

\paragraph{Coverage.} How much of the model does the equality test actually
reach? The FeynRules export covers a finite set of vertex arities, so some
operators may be structurally unreachable; see \Cref{sec:validation-coverage}.

Underneath all three sits an ordinary unit and regression suite for the
reusable engine: index validation, covariant expansion, field-strength
expansion, field transformations, fermion signs, tensor canonicalisation, and
selected benchmark Feynman rules.

Model-level comparisons live beside the concrete model packages under
\path{models/}, not inside the generic contraction engine: the engine
computes rules, and each model package owns the conventions and the
external-reference adapter needed to compare its own rules. The
format-independent comparison machinery itself is shared across models
(\Cref{sec:validation-generality}).

% --------------------------------------------------------------
\section{What ``equal'' means: the canonical form}
\label{sec:validation-canonical}
% --------------------------------------------------------------

Both sides of a comparison start as text: a FeynRules row is a
Mathematica-syntax string, such as \texttt{TensDot[Ga[mu],ProjM]}, and a
FeynPy row is produced by \texttt{feynman\_rule(...)} as described in
Chapter~\ref{ch:internals}. Both are parsed into Symbolica tensor
expressions before anything else happens; index brackets become index
symbols, \texttt{ProjM}/\texttt{ProjP} become chiral projectors, and indexed
Wilson parameters such as \texttt{alphaEcll[f1,f2,f3,f4]} become function
symbols \texttt{alphaEcll(f1,f2,f3,f4)}.

Two pieces of terminology are used throughout the rest of this chapter. The
\emph{signature} of a row is the ordered list of external field names, for
example \texttt{B|dR|dRbar} for the vertex coupling the hypercharge boson to
a right-handed down quark and its conjugate. The \emph{coefficient head} is
the name of the Wilson parameter multiplying a term, for example
\texttt{alphaEcll}. Comparison is always done per coefficient head within a
matching signature, never on a full undifferentiated sum (see below).

Equality is \emph{not} decided by comparing these parsed expressions
directly, nor by any string comparison. Two tensor expressions describing the
same physics can differ in ways that have nothing to do with whether the
physics is right: they can use different names for a summed-over (dummy)
index, write commuting factors in a different order, or leave the same
result differently expanded. Both sides are therefore reduced to a
\emph{canonical tensor-monomial map} --- a dictionary from canonical tensor
term to exact rational coefficient --- by
\path{src/symbolic/tensor_canonicalization.py}, and two rows are declared
equal exactly when their maps are identical, entry for entry, coefficient for
coefficient.

The reduction has four steps: (1) live tensor heads (the symbol names, e.g.\
the metric $g^{\mu\nu}$ or the Levi-Civita tensor $\epsilon^{\mu\nu\rho\sigma}$)
are swapped for temporary heads carrying explicit symmetry metadata, so the
canonicaliser is told which tensors are symmetric or antisymmetric under
index exchange; (2) Symbolica's tensor canonicaliser runs with contracted
dummy indices grouped, so their names carry no information; (3) dummy names
are standardised and commuting factors are sorted into one fixed order; (4)
the result is read off as the coefficient map.

\paragraph{Example 1: dummy names and symmetry.} The two expressions
\begin{equation}
    g^{\mu\nu}\, p_\mu k_\nu
    \qquad \text{and} \qquad
    g^{\nu\mu}\, k_\nu p_\mu
    \label{eq:validation-metric-example}
\end{equation}
are the same number for every $p,k$, but as raw Symbolica expressions they
differ in index naming ($\mu \leftrightarrow \nu$) and in the metric's index
order. Step (1) tells the canonicaliser that $g^{\mu\nu}=g^{\nu\mu}$; step (2)
identifies the two contracted-index pairs; only with both is the
canonicaliser able to reduce both expressions to the same map entry. Without
step (1), $g^{\mu\nu}$ and $g^{\nu\mu}$ would be treated as unrelated tensors
and the two expressions would wrongly be reported as different.

\paragraph{Example 2: generator reordering.} FeynRules and FeynPy can write
two colour or weak generators multiplying the same current in either order,
depending on the order in which the source Lagrangians introduced them. The
two orders are related by
\begin{equation}
    [T^a, T^b] = i f^{abc} T^c ,
    \label{eq:validation-commutator}
\end{equation}
so a term proportional to $T^b T^a$ on one side and $T^a T^b$ on the other
are equal only after adding the commutator term of
\Cref{eq:validation-commutator}. The canonicaliser applies exactly this one
identity when it encounters an out-of-order generator product
(\texttt{\_normalize\_generator\_product\_order\_report} in code); it does
not search over reorderings in general. The structure-constant Jacobi
identity,
\begin{equation}
    f^{abe}f^{cde} - f^{ace}f^{bde} + f^{ade}f^{bce} = 0 ,
    \label{eq:validation-jacobi}
\end{equation}
and the $SU(2)$ pseudoreality relations for $\epsilon_{ij}$ acting on a weak
generator are used the same narrow way, only where the two sides actually
need them.

These four steps and three identities are the entire simplification budget.
The comparison does \emph{not} use equations of motion, integration by parts,
momentum conservation, Schouten identities, Fierz identities, or any
open-ended pattern matching: admitting those would make the test far weaker,
because an incorrect rule could then often be massaged into the reference.

Two further design choices keep the comparison from degenerating into a
tautology. First, it runs \emph{separately for each coefficient head}: both
sides are filtered to the terms carrying a given head before the maps are
compared, so a cancellation between two different operators cannot mask an
error in either. Flavour order and complex conjugation stay inside the scalar
coefficient rather than being normalised away, so a term whose flavour
arguments have been silently transposed, e.g.\ $\alpha(f_2,f_1)$ for
$\alpha(f_1,f_2)$, is not judged equal merely because the coefficient head
name matches. Second, raw coefficient-head occurrence counts are recorded
only as a diagnostic, never as the equality criterion, because two equivalent
expanded forms can legitimately have different raw counts: FeynPy prints the
dual field strength $\mathrm{Dual}[F] = \tfrac{1}{2}\epsilon\cdot F$ as two
separate antisymmetric terms, where FeynRules has already used $\epsilon$
antisymmetry to collapse them into one. On the current SMEFT2 reference,
$100$ of $182$ shared signatures have matching raw counts, and every
remaining delta carries a recorded, specific reason; none is unexplained.

% --------------------------------------------------------------
\section{Standard Model benchmark}
\label{sec:validation-sm}
% --------------------------------------------------------------

The Standard Model package in \path{models/SM/} is the first complete
vertical slice of the project: a gauge-basis declaration, field
transformations to the physical basis, and a comparison of the resulting
interaction vertices against the bundled FeynRules \texttt{SM.fr} export.

The result is $163/163$ exact symbolic matches for the exported nonzero
flavour-expanded tree-level vertices, split by sector in
\Cref{tab:validation-sm-sectors}.

\begin{table}[htbp]
\centering
\begin{tabular}{lr}
\toprule
Sector & Vertices \\
\midrule
Gauge self-interactions & 8 \\
Matter currents & 51 \\
Yukawa & 42 \\
Higgs and Goldstone & 38 \\
Ghost & 24 \\
\midrule
Total & 163 \\
\bottomrule
\end{tabular}
\caption{Sector split of the Standard Model FeynRules comparison. Every row
in every sector is a direct exact symbolic match.}
\label{tab:validation-sm-sectors}
\end{table}

The scope should be stated precisely: the comparison runs with ghosts
included and gauge fixing disabled, so it validates the gauge, matter,
Yukawa, Higgs/Goldstone and ghost sectors but \emph{not} the gauge-fixing
sector, and it makes no claim about parameter-card semantics, loop-level
output, restriction files, or external formats such as UFO.

A second, independent benchmark is the unbroken background-field Standard
Model in \path{models/UnbrokenSM_BFM/}, matching its own FeynRules reference
at $67/67$ ($42$ three-point and $25$ four-point vertices). This matters for
SMEFT2, which is formulated in the same unbroken basis: it shows that the
basis, field naming, and generator conventions are already validated
independently of the EFT operators.

% --------------------------------------------------------------
\section{SMEFT2: model fidelity}
\label{sec:validation-fidelity}
% --------------------------------------------------------------

The main validation target is the SMEFT2 Green-basis model in
\path{models/SMEFT2/SMEFT2.py}, compared against the EFT-only FeynRules
\texttt{Ltot} export at
\path{models/SMEFT2/reference/Ltot_SMEFT_FeynRules.json}. The local FeynPy
\texttt{Ltot} follows the same convention, EFT contribution only, with the
SM-plus-EFT combination available separately as \texttt{Lfull}.

Before any vertex is compared, the transcription itself was audited. All
$32$ EFT Lagrangian blocks of \texttt{SMEFT\_Green\_Bpreserving.fr} are
implemented, and \texttt{Ltot} is composed from exactly those $32$ blocks in
the same order. All $298$ active Wilson coefficients of the FeynRules source
are declared in the FeynPy model and used in at least one operator; two
further coefficients, \texttt{alphaEcdlq} and \texttt{alphaEcqeu}, appear in
the FeynRules file but not in FeynPy, and both are commented out in
FeynRules and used in no FeynRules operator, so omitting them is correct
rather than a gap. (This is a statement about transcription, not about
comparison coverage: of the $298$ declared coefficients, the vertex
comparison of \Cref{sec:validation-smeft2} reaches $295$, for the structural
reason given in \Cref{sec:validation-coverage}.) The Standard Model core is
taken from \texttt{UnbrokenSM\_BFM.fr} minus its background-field ghost and
gauge-fixing sectors, deliberately excluded because the comparison target is
the EFT-only \texttt{Ltot}.

The structural differences between the two sources are systematic
conventions, not per-operator adjustments, as summarised in
\Cref{tab:validation-conventions}. This matters for
\Cref{sec:validation-ec}: because the FeynPy Lagrangian contains no
per-operator hand-fitting, a residual disagreement is evidence of a
\emph{convention} difference rather than of a transcription error.

\begin{table}[htbp]
\centering
\begin{tabular}{lll}
\toprule
FeynRules & FeynPy & Nature \\
\midrule
\texttt{2 Ta[...]} & \texttt{weak\_t} & named constant, 62 uses \\
\texttt{sigmamunu} & \texttt{sigma\_term} & named helper, 38 uses \\
\texttt{Dual[FS]} & \texttt{\_dual\_fs} & named helper \\
\texttt{lag + HC[lag]} & explicit \texttt{.conj()} terms & expansion \\
\texttt{CC[...]} & \texttt{dirac\_charge\_conjugation} & expansion \\
\bottomrule
\end{tabular}
\caption{Systematic convention differences between the FeynRules source and
the FeynPy transcription. No comparison-tuning comments, fudge factors, or
per-operator sign corrections appear in the model file.}
\label{tab:validation-conventions}
\end{table}

% --------------------------------------------------------------
\section{SMEFT2: comparison result}
\label{sec:validation-smeft2}
% --------------------------------------------------------------

All Green-basis and evanescent sectors used by the bundled reference are
included in the compiled FeynPy \texttt{Ltot}. The final result is given in
\Cref{tab:smeft2-final-comparison}; per-row verdicts are tabulated in
\Cref{app:smeft2-rows}.

\begin{table}[htbp]
\centering
\begin{tabular}{lr}
\toprule
Quantity & Result \\
\midrule
FeynRules reference rows & 184 \\
Operator-content matches & $184/184$ \\
\midrule
Direct exact symbolic matches & $161/184$ \\
Matches modulo the global \texttt{Ec} convention & $15/184$ \\
Matches modulo pinned row-specific packaging & $8/184$ \\
Unresolved packaging rows & $0/184$ \\
\midrule
Exact symbolic unequal rows & $0$ \\
Exact symbolic error rows & $0$ \\
Bosonic canonical tensor-map matches & $32/32$ \\
Bosonic canonical coefficient sectors & $93/93$ \\
Chirality-validated projectors & $2998$ \\
\bottomrule
\end{tabular}
\caption{Final SMEFT2 comparison summary for the EFT-only Green-basis
reference. The three-way exact-symbolic split is graded as described in
\Cref{sec:validation-ec}.}
\label{tab:smeft2-final-comparison}
\end{table}

The three-way split is the substance of this result and is reported
deliberately rather than collapsed into one number. A \emph{direct} match
means same signature and identical canonical tensor map, no packaging
assumption. A match \emph{modulo the global \texttt{Ec} convention} means the
row agrees only after the one fixed charge-conjugation transform of
\Cref{sec:validation-ec}. A match \emph{modulo pinned packaging} means the
row additionally needs one explicitly tabulated, row-specific transform, with
partner signature, phase and duplicate-leg symmetry all fixed in advance;
this group has two distinct sources, separated in
\Cref{sec:validation-ec-pinned}.

An earlier version of this comparison reported $176$ direct matches, because
the $15$ rows that depend on the global \texttt{Ec} convention were counted
as needing nothing. That was not wrong about equality, but it hid an
assumption behind an unconditional-looking number; the current grading keeps
the assumption visible.

The bosonic figure of $32/32$ is a diagnostic specific to the bosonic sector
and is not a coverage ceiling: the remaining $152$ fermionic rows are proved
by the same canonical-map machinery and are reported in the exact-symbolic
column above it.

% --------------------------------------------------------------
\section{The evanescent charge-conjugation convention}
\label{sec:validation-ec}
% --------------------------------------------------------------

The single most consequential assumption in the fermionic sector concerns
the evanescent charge-conjugated four-fermion operators, i.e.\ every
coefficient whose name begins with \texttt{alphaEc}. Without it, all $21$
four-fermion \texttt{Ec} rows fail, so it is documented here in full.

% --------------------------------------------------------------
\subsection{The discrepancy}
% --------------------------------------------------------------

FeynRules writes these operators through the \texttt{CC[...]} macro. Taking
the four-lepton operator as the representative case,
\begin{verbatim}
alphaEcll[f1,f2,f3,f4] CC[LLbar[sp1,ii,f1]].LL[sp1,jj,f2]
                       .LLbar[sp2,jj,f3].CC[LL[sp2,ii,f4]]
\end{verbatim}
\texttt{ExpandIndices} resolves the conjugation, so the exported vertex
carries \emph{no} residual charge-conjugation matrix: labelling the four
external legs $1,2,3,4$ in the order they appear, the spinor flow runs
$1\!\to\!4$ and $2\!\to\!3$ through ordinary spinor-metric and gamma chains.

FeynPy instead keeps the charge-conjugation matrix explicit and pairs
\emph{adjacent} legs,
\begin{equation}
    \bar{l}_L(1)\, C\, l_L(2)\;
    \bar{l}_L(3)\, C\, l_L(4) ,
    \label{eq:validation-ec-feynpy}
\end{equation}
that is, flow $(1,2)$ and $(3,4)$ with two explicit $C$ factors, where the
number in parentheses is the external-leg label.

% --------------------------------------------------------------
\subsection{The transform}
% --------------------------------------------------------------

The two forms describe the same operator, and the fix is now visible by
inspection: FeynPy's adjacent pairing $(1,2),(3,4)$ must be turned into
FeynRules's pairing $(1,4),(2,3)$, which is exactly the \emph{crossed}
re-pairing of the four legs. Eliminating the two $C$ factors this way picks
up one overall sign, from the antisymmetry of the charge-conjugation matrix,
\begin{equation}
    C^{T} = -C ,
    \label{eq:validation-c-antisymmetry}
\end{equation}
together with the anticommutation needed to reorder the fermion fields into
the crossed pairing.

Both the pairing and the sign are global constants, fixed once for every
\texttt{alphaEc} head rather than chosen per row: in code,
\texttt{\_EC\_CC\_CONVENTION\_MODE = "crossed"} and
\texttt{\_EC\_CC\_CONVENTION\_PHASE = -1}.

% --------------------------------------------------------------
\subsection{Why this is a derivation and not a fit}
\label{sec:validation-ec-overdetermination}
% --------------------------------------------------------------

Could a convention that rescues an entire sector just be a parameter tuned
until the answer came out right? The data rule this out: there are only four
possible (pairing, sign) combinations, and exactly one reproduces the
reference, uniformly, as shown in
\Cref{tab:validation-ec-overdetermination}.

\begin{table}[htbp]
\centering
\begin{tabular}{llr}
\toprule
Pairing mode & Phase & Rows reproduced (of 21) \\
\midrule
crossed & $-1$ & 15 \\
crossed & $+1$ & 0 \\
direct  & $-1$ & 3 \\
direct  & $+1$ & 2 \\
\bottomrule
\end{tabular}
\caption{Only one of the four possible conventions reproduces the FeynRules
reference. Under the accepted convention (crossed, phase $-1$) the transform
alone settles $15$ of the $21$ four-fermion \texttt{Ec} rows; the other $6$
additionally need the pinned partner transforms of
\Cref{sec:validation-ec-pinned}. The small counts under the three rejected
conventions are rows the transform does not apply to at all, not partial
successes.}
\label{tab:validation-ec-overdetermination}
\end{table}

The three rejected conventions fail on essentially every row; the contrast
is stark, not marginal. A per-row fit would have $4^{21}$ degrees of freedom,
whereas what is used is one global choice out of four, three of which the
data exclude. Removing the transform entirely fails all $21$ rows, which is
what makes it load-bearing and worth stating explicitly rather than leaving
it implicit.

% --------------------------------------------------------------
\subsection{Residual row-specific packaging}
\label{sec:validation-ec-pinned}
% --------------------------------------------------------------

Eight rows need a row-specific transform beyond what is in
\Cref{sec:validation-ec-overdetermination}, but they come from two unrelated
sources and should not be merged.

Two are Weinberg-operator rows, unconnected to the global \texttt{Ec}
convention: they carry the Weinberg coefficient, not an \texttt{alphaEc}
head, so the crossed-pairing transform never applies to them at all.
FeynRules packages them as same-chirality charge-conjugated lepton
structures, while FeynPy stores the mixed $\bar{l}_L/l_L$ form with the
charge-conjugation tensor explicit. The accepted transform compares the
FeynRules row against the antisymmetrised local pair
$\mathrm{FeynPy}(\bar{l}_L, l_L) - \mathrm{FeynPy}(l_L, \bar{l}_L)$, with the
relative minus sign fixed by \Cref{eq:validation-c-antisymmetry}, not
searched.

The remaining six rows \emph{are} additional to the global convention: they
are the six of the $21$ four-fermion \texttt{Ec} rows that the crossed-pairing
transform alone does not settle. They use a rule table of $12$
entries (one per coefficient sector), where each entry fixes exactly one
partner signature, one phase, one duplicate-leg symmetry mode, and a
provenance string naming the source operator; nothing is searched at
acceptance time. Dedicated sidecar checks confirm the Weinberg comparison at
$2/2$ and the \texttt{Ec} partner checks at $12/12$, and both sidecars also
verify that the \emph{wrong} sign or the wrong combination does \emph{not}
match, so they are evidence rather than restatement.

% --------------------------------------------------------------
\section{Chirality}
\label{sec:validation-chirality}
% --------------------------------------------------------------

The SMEFT2 parser discards the \texttt{ProjM}/\texttt{ProjP} chirality
projectors that FeynRules prints, because in the unbroken basis every matter
field is already chiral and the projector is redundant with the field
itself. That redundancy is a reasonable expectation, but not something to
assume silently: if a projector ever disagreed with the fields it connects,
dropping it would hide a genuine discrepancy. So the agreement is verified,
not assumed.

The underlying rule is elementary. Conjugation flips chirality, so a
left-handed field has a right-handed conjugate,
\begin{equation}
    \overline{P_X \psi} = \bar\psi\, P_{\bar X} ,
    \label{eq:validation-conjugate-chirality}
\end{equation}
where $\bar X$ is the chirality opposite $X$. Each gamma matrix anticommutes
with $\gamma_5$ and so swaps a projector for its opposite,
$\gamma^\mu P_X = P_{\bar X}\gamma^\mu$. Pushing the left projector of a
bilinear $\bar\psi_X\,\Gamma^{(n)}\,\psi_Y$ rightwards through its $n$ gamma
matrices gives
\begin{equation}
    \bar\psi_X \, \Gamma^{(n)} \, \psi_Y
    = \bar\psi \, \Gamma^{(n)}
    \begin{cases}
        P_{\bar X} P_Y \, \psi , & n \text{ even},\\
        P_{X} P_Y \, \psi , & n \text{ odd}.
    \end{cases}
    \label{eq:validation-chirality}
\end{equation}
Since $P_A P_B = \delta_{AB} P_B$, the chain survives only for $Y \neq X$
(even $n$) or $Y = X$ (odd $n$), and whenever it survives the one remaining
projector is $P_Y$. That prediction is exact and checkable against every
printed projector.

A chain whose two endpoints have the same ``barredness'' runs through a
charge-conjugation matrix instead of the spinor metric. Because $C$ itself
commutes with chirality, this flips the parity condition on $n$, and for a
barred/barred chain it also flips which chirality survives; both cases are
handled explicitly by the same rule.

The rule predicts every projector in the reference export with no
exceptions: all $2998$ of them, $2806$ ordinary chains and $192$
charge-conjugation chains. Any projector that contradicts the rule, or any
chain that must vanish identically for chirality reasons, raises an error
and the row is reported as failed. The check is model independent and lives
in the shared comparison layer, so other models can reuse it directly.

% --------------------------------------------------------------
\section{Mutation sensitivity}
\label{sec:validation-sensitivity}
% --------------------------------------------------------------

Everything above is a comparison that \emph{succeeded}, which on its own is
weak evidence: a comparison that had silently degenerated into a tautology
would also report a perfect score. The test suites therefore include
deliberate-corruption tests, which inject a physically meaningful error into
the reference and assert that the comparison notices.

For SMEFT2 the injected corruptions are: a rescaled Wilson coefficient (in
both the two-fermion and four-fermion sectors), a global sign flip,
chirality flips in both directions, transposed flavour slots, and a removed
imaginary unit. None is reported as agreement, and the same suite asserts
that the unmutated rows \emph{are} accepted, so it cannot pass vacuously by
failing everything.

The shared machinery carries an analogous suite under \path{models/SM/},
with six corruptions: a chirality flip, a wrongly conjugated CKM element, a
reversed colour generator, a removed imaginary unit, a wrong derivative
momentum, and a global sign flip. All six are detected.

Two structural facts reinforce the same point: no coefficient head is
exempted from comparison (the exemption set is empty), and the exception
tables are small and enumerated rather than open-ended. The raw
head-count deltas are explained by only four reason categories
(\texttt{DUAL\_FS\_ANTISYMMETRY}, \texttt{DUMMY\_LORENTZ\_MERGE},
\texttt{EXACT\_SYMBOLIC\_CANONICAL\_EQUIVALENCE},
\texttt{PINNED\_CC\_CANONICAL\_EQUIVALENCE}), and the pinned \texttt{Ec}
partner packaging is governed by exactly $12$ coefficient-sector rules, each
with a recorded provenance string.

% --------------------------------------------------------------
\section{Coverage and limits}
\label{sec:validation-coverage}
% --------------------------------------------------------------

The SMEFT2 model declares $303$ parameters: $5$ Standard Model parameters and
$298$ Wilson coefficients. Of these, $295$ (all Wilson coefficients) appear
in at least one validated reference vertex; the eight that do not split into
two groups.

Five are the Standard Model parameters themselves, $\lambda$, $\mu_H$ and the
three Yukawa matrices, correctly absent from an EFT-only \texttt{Ltot}.

The remaining three, \texttt{alphaKB}, \texttt{alphaR2B} and
\texttt{alphaOmuH2}, are genuine EFT coefficients the comparison cannot
reach, for a structural rather than accidental reason: the reference export
covers vertices of arity three to six, and these three coefficients multiply
operators that generate only two-point vertices (two from the abelian field
strength, whose non-abelian self-interactions are absent, and one from the
Higgs mass term). They are reported as unvalidated rather than counted as
covered, and a regression test pins this list so the gap cannot widen
unnoticed.

This is the complete coverage gap. Every other coefficient that contributes
to a two-point vertex also appears in a validated vertex of three or more
legs, where gauge invariance ties the two-point normalisation to the
higher-point ones. Closing the remaining gap would require exporting
two-point vertices from FeynRules, a straightforward extension of the
reference rather than a change to the comparison.

Two further limits: the comparison is tree-level and concerns vertex tensor
structure only, saying nothing about parameter-card semantics, loop-level
output, or external export formats; and the Standard Model benchmark
excludes the gauge-fixing sector, as noted in \Cref{sec:validation-sm}.

% --------------------------------------------------------------
\section{Generality of the comparison machinery}
\label{sec:validation-generality}
% --------------------------------------------------------------

The comparison is not SMEFT2-specific infrastructure. Its model-agnostic core
lives in \path{src/feynrules/comparison.py} and is shared by three models,
summarised in \Cref{tab:validation-generality}.

\begin{table}[htbp]
\centering
\begin{tabular}{llr}
\toprule
Model & Result & Adapter size \\
\midrule
\path{models/SM} & $163/163$ & 205 lines \\
\path{models/UnbrokenSM_BFM} & $67/67$ & 390 lines \\
\path{models/SMEFT2} & $184/184$ & 5530 lines \\
\bottomrule
\end{tabular}
\caption{The same comparison core serves three models. The Standard Model is
the reference integration and needs only a thin adapter; SMEFT2 is far larger
because it also carries the Wilson-coefficient sector logic, the
charge-conjugation packaging rules, and the sidecar checks.}
\label{tab:validation-generality}
\end{table}

Applying the machinery to a further model needs a parser adapter for that
model's export syntax, a field-name map, sector routing, and any convention
substitutions (CKM handling, diagonal Yukawa naming); models with Wilson
coefficients additionally need coefficient-head filtering. The machinery is
general in substance, though not plug-and-play.

SMEFT2 keeps its own parser because its export uses syntax the Standard
Model parsers do not cover: indexed Wilson coefficients, slashed momenta
inside spinor chains, two-index weak $\epsilon$ tensors, and generic index
deltas resolved by index prefix. Everything genuinely model independent is
shared rather than duplicated: bracket and comma scanning, FeynRules label
normalisation, the bosonic comparator, the canonical coefficient maps, and
the chirality validator of \Cref{sec:validation-chirality}.

% --------------------------------------------------------------
\section{Reproducibility}
\label{sec:validation-reproducibility}
% --------------------------------------------------------------

From the repository root, the full regression suite is
\begin{minted}{bash}
.venv/bin/python -m pytest -q
\end{minted}
which reports $519$ passing tests on the final checked state.

The accepted SMEFT2 gate is
\begin{minted}{bash}
.venv/bin/python -m models.SMEFT2.comparison --check --allow-cc-packaging
\end{minted}
which accepts the documented global \texttt{Ec} convention and the pinned
packaging rows, and fails on any unresolved packaging row, operator-content
miss, symbolic inequality, symbolic error, or unexplained FeynPy-only
signature.

The stricter audit gate omits the flag,
\begin{minted}{bash}
.venv/bin/python -m models.SMEFT2.comparison --check
\end{minted}
and requires every row to be a direct match with no packaging assumption at
all. It fails on the current state by construction, since $23$ rows are
documented packaging matches; its purpose is to detect whether that set has
changed, not to serve as the acceptance criterion.

The mutation-sensitivity suite of \Cref{sec:validation-sensitivity} runs
separately with
\begin{minted}{bash}
.venv/bin/python -m pytest \
  models/SMEFT2/tests/test_smeft2_comparison_sensitivity.py -q
\end{minted}

The human-readable report is regenerated with
\begin{minted}{bash}
.venv/bin/python -m models.SMEFT2.comparison
\end{minted}

Main reproducibility artifacts:
\begin{itemize}
    \item \path{models/SMEFT2/VALIDATION_ANALYSIS.md}, the maintained analysis
    of what the comparison proves and where its assumptions lie;
    \item \path{models/SMEFT2/COMPARISON.md}, the generated human-readable
    summary;
    \item \path{models/SMEFT2/COMPARISON_METHOD.md}, the maintained method
    report;
    \item \path{models/SMEFT2/comparison/artifacts/vertex_comparison_report.json},
    the per-signature comparison rows, from which \Cref{app:smeft2-rows} is
    generated;
    \item \path{models/SMEFT2/comparison/artifacts/feynpy_vertices.json}, the
    regenerated local FeynPy vertex rules.
\end{itemize}

% --------------------------------------------------------------
\section{Summary}
\label{sec:validation-summary}
% --------------------------------------------------------------

\paragraph{Established.} The FeynPy SMEFT2 Feynman rules agree with
FeynRules as tensor expressions, sector by sector in the Wilson coefficients,
for all $184$ exported vertices of arity three to six, covering $295$ of the
model's $298$ Wilson coefficients. Agreement is exact rational equality of
canonical tensor-monomial maps, is demonstrably sensitive to injected errors,
and rests on no unenumerated exceptions. The Standard Model and unbroken
background-field benchmarks match at $163/163$ and $67/67$.

\paragraph{Assumed, and documented.} Charge-conjugation packaging affects
$23$ of the $184$ rows, in three disjoint groups: $15$ rows settled by the
single global convention relating the FeynRules \texttt{CC[...]} flow to the
explicit $C$ factors of the FeynPy source; $2$ Weinberg rows using
antisymmetric $C$ packaging, independent of that convention; and $6$
evanescent partner rows needing both the global convention and a pinned
transform from a table of $12$ coefficient-sector rules. All three are
reported under distinct statuses rather than folded into the direct-match
count.

\paragraph{Outside scope.} The two-point sector, hence \texttt{alphaKB},
\texttt{alphaR2B} and \texttt{alphaOmuH2}; the Standard Model gauge-fixing
sector; and everything beyond tree-level vertex tensor structure.
